\documentclass[conference]{IEEEtran}
\IEEEoverridecommandlockouts
\usepackage{cite}
\usepackage{amsmath,amssymb,amsfonts}
\usepackage{algorithmic}
\usepackage{graphicx}
\usepackage{textcomp}
\usepackage{xcolor}
\usepackage{booktabs}
\usepackage{xeCJK}
\usepackage{hyperref}

\setCJKmainfont{PMingLiU}

\def\BibTeX{{\rm B\kern-.05em{\sc i\kern-.025em b}\kern-.08em
    T\kern-.1667em\lower.7ex\hbox{E}\kern-.125emX}}

\begin{document}

\title{DROS 4-Layer Defense-in-Depth Architecture for Autonomous AI Workloads\\
{\footnotesize \textnormal{面向 Agentic Web 時代的自主 AI 工作負載四層縱深防衛與執行治理架構}}
}

\author{\IEEEauthorblockN{陳俊誠 (Chun-Cheng "Jimmy" Chen)}
\IEEEauthorblockA{\textit{頂天立地股份有限公司 (Top-Celestial Company Ltd.)} \\
台北市, 台灣 \\
jimmychen@dr-os.io}
}

\maketitle
\thispagestyle{plain}
\pagestyle{plain}

\begin{abstract}
進入 2026 年，具備多步驟工具調用與自主推理能力之 AI Agent 已經深入企業金融合規、供應鏈調度與核心基礎設施等關鍵領域。然而，現有安全機制面臨嚴峻的架構邊界：應用層提示詞防火牆（如 NVIDIA NeMo Guardrails）本質上仍為機率性機制，極易遭間接提示詞注入（IPI）攻破；而核心層 OS 沙箱（如 eBPF、Seccomp）則面臨「上下文盲區（Context-Blindness）」，無法將使用者空間的 Agent 角色映射至底層系統呼叫。為彌補此「歸因缺口」，本文提出 \textbf{DROS 四層縱深防衛架構}，專為 Agentic Web 時代打造的零信任運行期執行控制面。本架構將控制面開通（經由 OpenAI Terraform Provider 與 OpenShip）與運行期物理執行防衛解耦，透過在二進位 C-ABI / FFI 邊界強制執行不可變之 $O(1)$ Capability Bitmap，達成亞微秒級決策延遲（\textbf{中位數 26.21 $\mu$s，熔斷延遲 $<$500 ns}）。配合三階 PKI 憑證中心發行具簽章之 \texttt{DrosIdentityToken (DIT)}，DROS 提供具 Ed25519 數位簽章之法庭級不可否認審計日誌，完全符合歐盟《EU AI Act》Sec. 50 法規要求。24 小時連續實證評測顯示，本架構對零日提示詞注入、目標挾持及跨企業供應鏈攻擊達成 100\% 確定性阻斷。
\end{abstract}

\begin{IEEEkeywords}
AI Agent 安全, 運行期執行治理, C-ABI 邊界強制阻斷, 零信任架構, 公鑰基礎設施 (PKI), 間接提示詞注入 (IPI), EU AI Act 合規性.
\end{IEEEkeywords}

\section{引言與問題定義}
從傳統被動對話式大語言模型演進至具備自主工具呼叫能力之 AI Agent，已徹底重塑企業的資安受攻擊面。現代自主 Agent 被賦予 API Key、OAuth 授權及資料庫交易權限。因此，當 Agent 遭挾持時，攻擊者係從授權邊界「內部」發動攻擊。

\subsection{身分與網路邊界防禦之失效}
傳統網頁應用程式防火牆 (WAF)、端點偵測回應 (EDR) 與身分存取管理 (IAM) 系統均建立於「攻擊者未具備合法憑證」之假設上。但在 Agentic Web 時代，遭挾持的 Agent \textit{即為}持證合法實體。

\subsection{語意與核心之悖論 (Semantic-Kernel Paradox)}
如最新資安文獻（如 AgentSight, IEEE S\&P 2026）所指出，AI Agent 安全面臨根本性悖論：
\begin{enumerate}
    \item \textbf{語意防火牆（高語意，零確定性）：} 提示詞過濾器可檢查自然意圖，但無法提供抗混淆之數學確定性。
    \item \textbf{核心沙箱（高確定性，零語意）：} eBPF 或 Seccomp 等 OS 機制可執行嚴格二進位規則，但缺乏使用者空間上下文，無法辨識系統呼叫屬於何種角色。
\end{enumerate}

\begin{verbatim}
┌─────────────────────────────────────────────────────────────────────────────┐
│ 1. Control Plane & GitOps Provisioning (Policy-as-Code)                     │
│    • OpenAI Terraform Provider -> Provision Projects, IAM, Keys & Rate Limits│
│    • OpenShip Engine           -> Orchestrate Multi-Enterprise Containers   │
├─────────────────────────────────────────────────────────────────────────────┤
│ 2. DROS Runtime Execution Defense (L4 - C-ABI Boundary Enforcement)          │
│    • 3-Tier PKI Certificate Chain -> DrosIdentityToken (DIT) Binding       │
│    • DROS GuardVM (PEP/PDP)       -> Sub-microsecond <500ns Binary Panic    │
└─────────────────────────────────────────────────────────────────────────────┘
\end{verbatim}

\section{威脅模型：Agent 運行期攻擊向量 (AAV-2026)}
我們將針對自主 AI Agent 的運行期攻擊向量定義為 \textbf{Agentic Attack Vectors (AAV-2026)}：

\begin{table}[htbp]
\caption{Agent 運行期攻擊向量分類表}
\begin{center}
\begin{tabular}{|p{2.2cm}|p{1.8cm}|p{3.5cm}|}
\hline
\textbf{攻擊向量} & \textbf{MITRE ATLAS} & \textbf{作用機制與影響} \\ \hline
\textbf{間接提示詞注入 (IPI)} & AML.T0051 & 攻擊者將惡意指令植入外部資料，劫持工具呼叫流程。 \\ \hline
\textbf{目標劫持 (Goal Hijacking)} & AML.T0054 & 上下文視窗污染改變長期目標，引發未授權動作鏈。 \\ \hline
\textbf{特權函數越權} & AML.T0053 & 遭挾持 Agent 利用合法憑證呼叫超出角色權限之高特權 API。 \\ \hline
\textbf{供應鏈傳染} & AML.T0010 & 受污染之外鄰套件/模型挾持擷取 Agent，進而轉向攻擊內部資產。 \\ \hline
\end{tabular}
\end{center}
\end{table}

\section{DROS 四層縱深防衛架構}
DROS 建立四層縱深防衛模型，將機率性過濾器與確定性執行門檻清晰解耦：

\begin{verbatim}
┌──────────────────────────────────────────────────────────────────────┐
│ L1: Detective Intelligence Layer  │ Semantic Prompt Filtering (~90%)  │
├───────────────────────────────────┼──────────────────────────────────┤
│ L2: Zero Trust Mesh & PKI Layer   │ 3-Tier CA (Root->AIA->BEC) + DIT  │
├───────────────────────────────────┼──────────────────────────────────┤
│ L3: Task Orchestration Layer      │ Multi-Agent Swarm ABAC Isolation │
├───────────────────────────────────┼──────────────────────────────────┤
│ ★ L4: C-ABI Physical Enforcement  │ <500ns O(1) Binary Panic Gate    │
└───────────────────────────────────┴──────────────────────────────────┘
\end{verbatim}

\subsection{L1: 偵測情報層}
L1 透過語意分析清洗自然語言輸入，阻斷已知提示詞注入與越獄樣板。因 L1 為機率性機制，故定位為前置過濾。

\subsection{L2: 零信任網格與 PKI 身分層}
為消除上下文盲區，L2 引入三階 PKI 憑證體系（Root CA $\to$ AIA $\to$ BEC Leaf），將 Agent 身分、角色與 Skill 授權綁定於 signed \textbf{DrosIdentityToken (DIT)} 中。

\subsection{L3: 任務調度層}
L3 經由 \texttt{agent\_manifest.yaml} 對多 Agent Swarm 強制執行屬性型存取控制 (ABAC)，將爆炸半徑隔離於預先核可之圖譜拓撲內。

\subsection{L4: C-ABI 物理硬熔斷層}
L4 為 DROS 之核心創新。權限於初始化時預編譯為不可變之 Capability Bitmap。當工具呼叫觸發時，GuardVM 執行 $O(1)$ 位元與運算：
\begin{equation}
\text{Decision} = \text{Capability\_Bitmap}[\text{Role\_ID}] \ \text{\&} \ \text{Requested\_Tool\_Bit}
\end{equation}
若位元為 $0$，執行將於 C-ABI 二進位邊界以 \textbf{$<$500 ns 熔斷延遲} 實體硬阻斷。

\section{聯邦式 B2B 跨企業供應鏈防衛}
當自主 Agent 跨越企業邊界互動時（如 Corp-Alpha 存取 Corp-Beta 資料），DROS 將 L2 轉化為跨域身分指紋門檻，提供細胞級爆炸半徑收斂與微秒級憑證廢止 (CRL)。

\section{實驗評測與基準結果}

\subsection{開源評測靶場}
所有實證評測均於 \textbf{DROS-VEP Lite} 開源容器化環境執行。測試腳本已開源於 \texttt{github.com/Top-Celestial-Company-Ltd/DROS-VEP-lite}。

\subsection{反事實對照組基準評測}

\begin{table}[htbp]
\caption{反事實對照組實驗結果 (無防衛 vs. 啟用 DROS Guard L4)}
\begin{center}
\begin{tabular}{|p{2.3cm}|p{2.2cm}|p{2.2cm}|p{1.0cm}|}
\hline
\textbf{評測劇本} & \textbf{對照組 (無防禦)} & \textbf{實驗組 (DROS L4)} & \textbf{攔截延遲} \\ \hline
\textbf{ATS-001 (EP1 資料)} & ❌ 100\% 外洩 & ✅ 100\% 阻斷 & 25.8 $\mu$s \\ \hline
\textbf{ATS-002 (EP2 密碼)} & ❌ 100\% 遭竊 & ✅ 100\% 阻斷 & 26.1 $\mu$s \\ \hline
\textbf{ATS-003 (EP3 部署)} & ❌ 100\% 部署 & ✅ 100\% 阻斷 & 25.5 $\mu$s \\ \hline
\textbf{ATS-004 (EP4 供應鏈)} & ❌ 100\% 挾持 & ✅ 100\% 阻斷 & 26.4 $\mu$s \\ \hline
\end{tabular}
\end{center}
\end{table}

\subsection{24 小時實體浸泡測試微觀數據}

\begin{table}[htbp]
\caption{24 小時連續浸泡測試遠測數據}
\begin{center}
\begin{tabular}{|p{3.5cm}|p{2.2cm}|p{2.0cm}|}
\hline
\textbf{評測指標} & \textbf{24小時實測數據} & \textbf{備註} \\ \hline
\textbf{總評測請求負載} & \textbf{160,611 次} & 連續 24.0 小時 \\ \hline
\textbf{策略評估延遲 (P50)} & \textbf{26.21 $\mu$s} & $\pm 0.34\ \mu\text{s}$ \\ \hline
\textbf{P99 策略延遲} & \textbf{242.69 $\mu$s} & 高並發峰值 \\ \hline
\textbf{C-ABI 實體熔斷延遲} & \textbf{$<$ 500 ns} & 二進位門檻 \\ \hline
\textbf{SPEC CPU2017 負載} & \textbf{$<$ 1.8\%} & 極低上下文切換 \\ \hline
\textbf{確定性阻斷成功率} & \textbf{100\% (0外洩)} & 攔截137,751次 \\ \hline
\textbf{24小時記憶體洩漏} & \textbf{0 Bytes} & 零堆積分配 \\ \hline
\end{tabular}
\end{center}
\end{table}

\section{相關文獻}
本研究延伸四大資安領域：系統呼叫攔截 (eBPF AgentSight)、LLM Agent 安全 (OWASP Top 10 for Agentic Applications)、零信任架構 (NIST SP 800-207) 與不可否認性審計基質 (EU AI Act Sec. 50)。

\section{結論與未來展望}
本文提出 DROS 四層縱深防衛架構，建立結合三階 PKI 身分驗證與亞微秒級 C-ABI 二進位強制阻斷之確定性運行期控制面。連續 24 小時實證評測證實本架構達成 100\% 確定性攔截。

\section*{致謝與 AI 協作宣告}
本研究之核心安全模型、C-ABI 邊界阻斷範式及美國專利主張（U.S. PPA No. 64/111,973）均由陳俊誠 (Jimmy Chen) 獨立提出與驗證。生成式 AI 工具僅嚴格用於排版與語法輔助。

\begin{thebibliography}{00}
\bibitem{b1} NIST SP 800-207, ``Zero Trust Architecture,'' NIST, 2020.
\bibitem{b2} OWASP Foundation, ``OWASP Top 10 for LLM Applications v1.1,'' 2023.
\bibitem{b3} MITRE Corporation, ``MITRE ATLAS,'' 2024.
\bibitem{b4} X. Zhang et al., ``AgentSight: eBPF-Powered Tracing for LLM Agents,'' arXiv:2408.01234, 2024.
\bibitem{b5} C. C. Chen, ``Runtime Attribution Framework,'' Zenodo, DOI: 10.5281/zenodo.20823163, 2026.
\bibitem{b6} C. C. Chen, ``DROS-PGM,'' U.S. Patent Application No. 64/111,973, 2026.
\bibitem{b7} European Parliament and Council, ``Regulation (EU) 2024/1689 (EU AI Act),'' 2024.
\end{thebibliography}

\end{document}
